%https://de.overleaf.com/learn/latex/Glossaries infos über Erstellung glossareintraäge und aufrufe dieser
%WICHTIG! wenn Änderungen an den Glossareinträgen vorgenommen werden,im Terminal hintereinander das hier iengeben:
%erstmal .glo datei lsöchen, dann:
%pdflatex Seminarfacharbeit.tex
%makeglossaries Seminarfacharbeit
%pdflatex Seminarfacharbeit.tex
%nur dann funktioniert das Glossar auch in der PDF-Datei
\begin{comment}
    zum kopieren:
    \newglossaryentry{}
    {}
\end{comment}

\newglossaryentry{silber-silberchlorid-elektrode}
    {
    name=Silber-Silberchlorid-Elektrode,
    text={\textit{Silber-Silberchlorid-Elektrode}},
    first={\textit{Silber-Silberchlorid-Elektrode}},
    description={Die Silber-Silberchlorid-Elektrode ist eine häufig verwendete Elektrode in der Elektroenzephalographie und besteht aus einem Silberdraht, der mit einer Schicht aus Silberchlorid überzogen ist}
    }

\newglossaryentry{berger-effekt}
{
    name=Berger-Effekt,
    text={\textit{Berger-Effekt}},
    first={\textit{Berger-Effekt}},
    description={Der Berger-Effekt bezeichnet eine Blockade der Alpha-Wellen im EEG und ist nachweisbar, in dem ein Proband, der wach und entspannt ist, die Augen öffnet. Die normalerweise bei geschlossenen Augen sichtbaren Alpha-Wellen verschwinden und werden durch Beta-Wellen ersetzt. Beim Schließen der Augen kehrt der Alpha-Rhythmus zurück}
}

\newglossaryentry{brain-computer-interfaces}
{
    name=Brain-Computer-Interfaces,
    text={\textit{Brain-Computer-Interfaces}},
    first={\textit{Brain-Computer-Interfaces}},
    description={Brain-Computer-Interfaces (BCI) sind Systeme, die es ermöglichen, Informationen direkt aus dem Gehirn zu extrahieren und in Befehle für externe Geräte umzuwandeln}
}

\newglossaryentry{elektrolyt}
{
    name=Elektrolyt,
    text={\textit{Elektrolyt}},
    first={\textit{Elektrolyt}},
    description={Ein Elektrolyt ist eine Substanz, die in der Lage ist, elektrische Ladungen zu leiten, wenn sie als Lösung oder im geschmolzenem Zustand vorliegt}
}

\newglossaryentry{analog-digital-converter}
    {name=Analog-Digital-Converter,
    text={\textit{Analog-Digital-Converter}},
    first={\textit{Analog-Digital-Converter}},
    description={Ein Analog-Digital-Converter, auf deutsch Analog-Digital-Wandler, ist ein elektronisches Bauteil, das analoge Signale in einen digitalen Datenstrom, der dann weiterverarbeitet oder gespeichert werden kann}
}

\newglossaryentry{notch}{
    name=Notch-Filter,

    description={Filter, der gezielt eine bestimmte Frequenz (z.B. 50 Hz Netzrauschen) unterdrückt}
}

\newglossaryentry{ica}{
    name=ICA,
    description={Independent Component Analysis, mathematisches Verfahren zur Trennung überlagerter Signalquellen im EEG}
}

\newglossaryentry{bandpass}{
    name=Bandpassfilter,
    description={Filter, der nur Signale innerhalb eines bestimmten Frequenzbereichs durchlässt}
}

\newglossaryentry{cutoff}{
    name=Cut-off-Frequenz,
    description={Grenzfrequenz, ab der ein Filter Signale durchlässt oder unterdrückt}
}

\newglossaryentry{p300}{
    name=P300,
    description={ERP-Komponente, die etwa 300 ms nach einem seltenen oder bedeutsamen Stimulus auftritt}
}

\newglossaryentry{eog}{
    name=EOG,
    description={Elektrookulogramm, Messung von Augenbewegungen zur Identifikation von Artefakten im EEG}
}

\newglossaryentry{eventmarker}{
    name=Ereignismarker,
    description={Zeitliche Markierung im EEG-Signal zur Kennzeichnung eines Stimulus oder Ereignisses}
}

\newglossaryentry{oddball}{
    name=Oddball-Aufgabe,
    description={Experimentelles Paradigma, bei dem seltene Zielreize zwischen häufigen Standardreizen präsentiert werden, um ERPs zu untersuchen}
}
